\begin{definition}[LCM period for a finite block family]
    \label{def:lcm_period}
    \lean{MPSTensor.lcmPeriod}
    \leanok
    For a finite family of positive integer periods $(p_i)_{i \in I}$, their
    least common multiple $\lcm_{i \in I}(p_i)$ defines the shared blocking
    period.
\end{definition}

\begin{definition}[Common-period blocked family]
    \label{def:common_period_blocking}
    \lean{MPSTensor.commonPeriodBlocking}
    \leanok
    \uses{def:lcm_period, def:blocked_tensor}
    For a finite family of tensors $(B_i)_{i \in I}$ with periods
    $(p_i)_{i \in I}$, the \emph{common-period blocked family} consists of the
    blocked tensors $B_i^{[L]}$ with $L = \lcm_{j \in I}(p_j)$.
\end{definition}

\begin{theorem}[Common-period blocking preserves transfer-map primitivity]
    \label{thm:common_period_blocking_primitive}
    \lean{MPSTensor.isPrimitive_transferMap_commonPeriodBlocking}
    \leanok
    \uses{def:common_period_blocking, thm:iterated_blocking_transfer, thm:primitive_blocking_divisible}
    If each tensor $B_i$ has a primitive transfer map after blocking by its
    period $p_i \ge 1$, then each blocked tensor $B_i^{[L]}$ with
    $L = \lcm_j(p_j)$ also has a primitive transfer map.
\end{theorem}

\begin{proof}\leanok
    Because each $p_i$ divides $L = \lcm_j(p_j)$,
    Theorem~\ref{thm:primitive_blocking_divisible} implies that the transfer map
    $\E_{B_i^{[L]}}$ is primitive for every $i \in I$.
\end{proof}
